% =====================================================================================================
%  قالب پایان‌نامه‌ی فارسی (XeLaTeX + xepersian، قلم یاس)
%  کامپایلر Overleaf: XeLaTeX     |     ساخت محلی: scripts/build.sh main.tex
%  تصاویر با \figbox{name.png}{width} درج می‌شوند؛ اگر فایل نباشد کادر جای‌خالی چاپ می‌شود.
% =====================================================================================================
\documentclass[12pt,a4paper]{report}

\usepackage{amsmath,amssymb}
\usepackage{graphicx}
\usepackage[a4paper,top=2.8cm,bottom=2.6cm,right=2.6cm,left=2.4cm,headheight=16pt]{geometry}
\usepackage{setspace}
\usepackage{array}
\usepackage{booktabs}
\usepackage{enumitem}
\usepackage{fancyhdr}
\usepackage[colorlinks=true,linkcolor=black,citecolor=blue,urlcolor=blue]{hyperref}
\usepackage{xepersian}   % باید آخرین بسته باشد

% ---------------------------------------- قلم یاس ----------------------------------------------------
% به ترتیب امتحان می‌شود: فایل قلم بارگذاری‌شده در پروژه، سپس قلم‌های نصب‌شده با نام‌های رایج «یاس».
% Scale=1.12 متن را حدود ۱۲ درصد بزرگ‌تر از نسخه‌ی قبلی می‌کند.
% اگر در پروژه‌ی قبلی‌تان خطی مثل \settextfont{...} داشتید، می‌توانید کل این بلوک را با همان خط
% (با افزودن Scale=1.12) جایگزین کنید.
\IfFontExistsTF{BYas.ttf}{%
  \IfFontExistsTF{BYasBd.ttf}%
    {\settextfont[Scale=1.12,BoldFont=BYasBd.ttf]{BYas.ttf}}%
    {\IfFontExistsTF{BYasBold.ttf}%
      {\settextfont[Scale=1.12,BoldFont=BYasBold.ttf]{BYas.ttf}}%
      {\settextfont[Scale=1.12]{BYas.ttf}}}}%
{\IfFontExistsTF{BYas}{\settextfont[Scale=1.12]{BYas}}%
{\IfFontExistsTF{B Yas}{\settextfont[Scale=1.12]{B Yas}}%
{\IfFontExistsTF{XB Yas}{\settextfont[Scale=1.12]{XB Yas}}%
{\IfFontExistsTF{Yas}{\settextfont[Scale=1.12]{Yas}}%
{\IfFontExistsTF{XB Niloofar}{\settextfont[Scale=1.12]{XB Niloofar}}%
{\IfFontExistsTF{Vazirmatn}{\settextfont[Scale=1.05]{Vazirmatn}}%
{\IfFontExistsTF{Nazli}{\settextfont[Scale=1.12]{Nazli}}%
{\settextfont[Scale=1.05]{Amiri}}}}}}}}}

% ---------------------------------------- فاصله‌ی خطوط --------------------------------------------
\setstretch{1.55}
\setlength{\parskip}{3pt}
\setlist{itemsep=4pt,topsep=4pt}
\setlength{\emergencystretch}{3em}

% ---------------------------------------- سرصفحه ---------------------------------------------------
\pagestyle{fancy}
\fancyhf{}
\fancyhead[R]{\small\leftmark}
\fancyfoot[C]{\thepage}
\renewcommand{\headrulewidth}{0.5pt}
\fancypagestyle{plain}{\fancyhf{}\fancyfoot[C]{\thepage}\renewcommand{\headrulewidth}{0pt}}

% ---------------------------------------- فرمان‌های کمکی ------------------------------------------
\newcommand{\sub}[1]{\mathrm{#1}}               % زیرنویس لاتین در فرمول‌ها (مثل rad، tot)
\newcommand{\Imag}{\operatorname{Im}}
\newcommand{\Real}{\operatorname{Re}}
% محل تصویر: اگر فایل وجود داشته باشد درج می‌شود، وگرنه کادر خالی با نام فایل
\newcommand{\figbox}[2]{%
  \IfFileExists{#1}{\includegraphics[width=#2]{#1}}{%
    \fbox{\parbox[c][\dimexpr#2*42/100\relax][c]{\dimexpr#2-2\fboxsep-2\fboxrule\relax}{\centering\textbf{محل قرارگیری نمودار لومریکال}\\[4pt]
    فایل تصویر را با نام \lr{\texttt{\detokenize{#1}}} بارگذاری کنید.}}}}

\begin{document}

% ======================================== صفحه‌ی عنوان ============================================
\begin{titlepage}
\begin{center}
\vspace*{1.5cm}
{\Large\textbf{دانشگاه ...}}\\[6pt]
{\large دانشکده‌ی ...}\\[4pt]
{\large گروه ...}\\[2.2cm]
\rule{\textwidth}{0.6pt}\\[10pt]
{\LARGE\textbf{عنوان پایان‌نامه}}\\[14pt]
{\large\textbf{زیرعنوان}}\\[10pt]
\rule{\textwidth}{0.6pt}\\[2.5cm]
\begin{minipage}[t]{0.48\textwidth}
\textbf{پژوهشگر:}\\ نام دانشجو\\[12pt]
\textbf{شماره‌ی دانشجویی:}\\ ۰۰۰۰۰۰۰۰۰۰
\end{minipage}\hfill
\begin{minipage}[t]{0.42\textwidth}
\textbf{استاد راهنما:}\\ نام استاد
\end{minipage}
\end{center}
\end{titlepage}

% ======================================== چکیده =====================================================
\chapter*{چکیده}
\addcontentsline{toc}{chapter}{چکیده}
\markboth{چکیده}{چکیده}
متن چکیده (بدون ارجاع). واژه‌های لاتین همیشه داخل \lr{Latin text}.

\vspace{6pt}
\noindent\textbf{واژگان کلیدی:} ...، ...، ....

\tableofcontents

% ======================================== فصل‌ها ======================================================
\chapter{مقدمه}
\section{بیان مسئله}
متن با ارجاع \cite{key1}. معادله:
\begin{equation}\label{eq:example}
F_P=\frac{P_{\sub{tot}}}{P_0}
\end{equation}
ارجاع به رابطه‌ی (\ref{eq:example})، شکل \ref{fig:example} و جدول \ref{tab:example}.

\begin{figure}[htbp]
\centering
\figbox{example.png}{0.85\textwidth}
\caption{توضیح شکل؛ واحدها به فارسی (نانومتر)، اعشار با «/» مثل ۵۴/۵.}
\label{fig:example}
\end{figure}

\begin{table}[htbp]
\centering
\caption{عنوان جدول (بالای جدول).}
\label{tab:example}
\renewcommand{\arraystretch}{1.3}
\resizebox{\textwidth}{!}{%
\begin{tabular}{ccc}
\toprule
ستون ۱ & ستون ۲ & ستون ۳ \\
\midrule
۳ & ۴۲۰۰ & حدود ۶ درصد \\
\bottomrule
\end{tabular}}
\end{table}

% ======================================== منابع ======================================================
\chapter*{منابع و مآخذ}
\addcontentsline{toc}{chapter}{منابع و مآخذ}
\markboth{منابع و مآخذ}{منابع و مآخذ}
\begin{latin}
\renewcommand{\bibname}{\rl{منابع و مآخذ}}
\renewcommand{\chapter}[2]{}   % عنوان لاتین تکراری حذف می‌شود
\begin{thebibliography}{99}
\bibitem{key1} A. Author, ``Title,'' \textit{Journal}, vol. 1, no. 1, pp. 1--10, 2020.
\end{thebibliography}
\end{latin}

\end{document}
